As in many applications in plasma physics, collisions dominate the thermal
particles, and therefore, the distribution function may be expanded about a
Maxwellian
\begin{equation}
f_{0}=f_{0M}+\widehat{f}_{0}
\end{equation}
and the self-collision operator $\mathcal{C}\left( f_{0}\right) \equiv 
\mathcal{C}\left( f_{0},f_{0}\right) $may be approximated by the linearized
operator 
\begin{equation}
\mathcal{C}\left( f_{0},f_{0}\right) \simeq \mathcal{C}_{lin}\left(
f_{0},f_{0}\right) \simeq \mathcal{C}\left( f_{0},f_{0M}\right) +\mathcal{C}%
\left( f_{0M},f_{0}\right) 
\end{equation}
where one have made use of the fact that $\mathcal{C}\left(
f_{0M},f_{0M}\right) =0$ while  $\widehat{f}_{0}^{2}$ terms are ignored. 

The usual approach is to express $\mathcal{C}\left( f_{0M},f_{0}\right) $ as
a sum of Legendre harmonics, 
\begin{equation}
\mathcal{C}(f_{0M},f_{0})=\sum_{l=0}^{\infty }\mathcal{C}\left(
f_{0M},f_{0,l}P_{l}\left( \xi \right) \right) 
\end{equation}
where
\begin{equation}
f_{0}\left( p,\xi ,\psi ,\theta \right) =\sum_{l=0}^{\infty %
}f_{0,l}P_{l}\left( \xi \right) 
\end{equation}
with
\begin{equation}
f_{0,l}\left( p,\xi ,\psi ,\theta \right) =\left( l+\frac{1}{2}\right)
\int_{-1}^{+1}f_{0}\left( p,\xi ,\psi ,\theta \right) P_{l}\left( \xi
\right) d\xi 
\end{equation}

Here, 
\begin{equation}
\mathcal{C}\left( f_{0},f_{0}\right) =\mathcal{C}_{ee}\left(
f_{0},f_{0}\right) +\mathcal{C}_{ei}\left( f_{0},f_{0}\right) 
\end{equation}
where $\mathcal{C}_{ei}$ and $\mathcal{C}_{ee}$ describe respectively the
electron-ion and electron-electron collision processes. Because of the large
difference of mass between electrons and ions, $\mathcal{C}%
_{ei}(f_{0M},f_{0})\simeq 0,$ while higher order terms must be kept for
electron-electron interactions, in order to conserve momentum transfer
during the collisions. For neary all current drive problems in tokamak
physics, the sum may be truncated up to the first term\footnote{%
Second order calculations have shown negligible modifications on the
numerical solution of the drift kinetic euqation for the current drive
problem.}, and therefore,
\begin{equation}
\mathcal{C}_{ee}\left( f_{0M},f_{0}\right) \simeq \mathcal{C}_{ee}\left(
f_{0M},\xi f_{0_{1}}\right) \simeq \xi \mathcal{I}\left(
f_{0M},f_{0,l=1}\right) 
\end{equation}
where
\begin{equation}
f_{0,l=1}\left( p,\xi ,\psi ,\theta \right) =\frac{3}{2}\int_{-1}^{+1}\xi
f_{0}\left( p,\xi ,\psi ,\theta \right) d\xi 
\end{equation}

As a consequence, the collision operator conserves number and momentum, but
not energy, so there is no need to introduce an energy loss term in the
equations. In its general form, the expression of $\mathcal{C}\left(
f_{0}\right) $is then
\begin{equation}
\mathcal{C}\left( f_{0}\right) \simeq \mathcal{C}_{ee}\left(
f_{0},f_{0M}\right) +\mathcal{C}_{ei}\left( f_{0},f_{0M}\right) +\mathcal{C}%
_{ee}\left( f_{0M},\xi f_{0,l=1}\right) 
\end{equation}
where the first two terms may be cast as the divergence of a flux as follows 
\begin{equation}
\mathcal{C}_{ee}(f_{0},f_{0M})+\mathcal{C}_{ei}(f_{0},f_{0M})=-\nabla _{%
\mathbf{p}}\cdot \mathbf{S}^{\text{C}}=-\frac{1}{p^{2}}\frac{\partial }{%
\partial p}\left( p^{2}S_{p}^{\text{C}}\right) +\frac{1}{p}\frac{\partial }{%
\partial \xi }\left( \sqrt{1-\xi ^{2}}S_{\xi }^{\text{C}}\right) 
\end{equation}
with the components 
\begin{eqnarray}
S_{p}^{\text{C}} &=&-\left[ A\left( p,\psi ,\theta \right) \frac{\partial %
f_{0}}{\partial p}+F\left( p,\psi ,\theta \right) f_{0}\right]  \\
S_{\xi }^{\text{C}} &=&\frac{B_{t}\left( p,\psi ,\theta \right) }{p}\sqrt{%
1-\xi ^{2}}\frac{\partial f_{0}}{\partial \xi }
\end{eqnarray}

Coefficients $A\left( p,\psi ,\theta \right) $ and  $F\left( p,\psi ,\theta
\right) $ describes respectively momentum diffusion and convection, while $%
B_{t}\left( p,\psi ,\theta \right) $corresponds to pitch-angle scattering.
There values as well as that of $\mathcal{I}\left( p,\psi ,\theta \right)
\equiv \mathcal{I}\left( f_{0M},f_{0,l=1\ss }\right) $depend of a model of
collision. Simplified expressions may be obtained for an isotropic or
Maxwellian backgrounds and also in the high velocity limit, when $v\gg
v_{th}.$ The latter case has been fruitfuly used for the current drive
problem involving very energetic or runaway electrons. Detailed expressions
of  $A\left( p,\psi ,\theta \right) ,$ $F\left( p,\psi ,\theta \right) $,  $%
B_{t}\left( p,\psi ,\theta \right) $and $\mathcal{I}\left( p,\psi ,\theta
\right) $are given in Appendix *** for this useful limit. 

In the code, the coefficients are deduced from the Beliaev-Budker collision
operator which is a fully relativistic formulation of the Landau expression
which is only valid for weakly energetic particles. Despite its complex
tensorial form, it has been successfuly expressed as pseudo-Rosenbluth
potentials using a Legendre harmonics expansion [B. Brams and C. F. Karney,
Phys. Fluids B1 (1989) 1355], a formulation that opens the possibility to
incorporate this operator in numerical calculations, since integrals may be
determined in a fast and accurate manner. As compared to methods used in
most of the kinetic codes dedicated to the current drive problem, which
usually use several approximate expressions for the collision operator only
valid in limited range of energy, the use of the  Beliaev-Budker collision
operator represents a great advantage. Indeed, the whole momentum dynamics
from non-relativistic to highly relativistic limit may be described by a
single operator, which strongly enhance simulation reliability, when both
limits come simultaneously into play,like for self-consistent bootstrap
current calculations in presence of RF waves. Exact expressions of  $A\left(
p,\psi ,\theta \right) $, $F\left( p,\psi ,\theta \right) $, $B_{t}\left(
p,\psi ,\theta \right) $and $\mathcal{I}\left( p,\psi ,\theta \right) $are
given in Appendix ****. 

Finally, it is important to point out that the drift kinetic equation is no
more a simple partial differential equation, but an integro-differential
equation. However, since the term  $\mathcal{C}_{ee}\left(
f_{0M},f_{0}\right) $introduce only a weak non-linearity, it may be
considered as an additional source tem in the equations.  

According to the conservative formulation, the collisional flux can be
divided into diffusive and convective contributions, which in tensorial form
gives 
\[
\mathbf{S}^{\text{C}}=-\Bbb{D}_{\mathbf{p}}^{\text{C}}\cdot \nabla _{p}f_{0}+%
\mathbf{F}_{\mathbf{p}}^{\text{C}}f_{0}
\]
with the diffusion tensor components 
\begin{eqnarray}
D_{pp}^{\text{C}} &=&A\left( p,\psi ,\theta \right)   \label{Dcoll} \\
D_{p\xi }^{\text{C}} &=&0 \\
D_{\xi p}^{\text{C}} &=&0 \\
D_{\xi \xi }^{\text{C}} &=&B_{t}\left( p,\psi ,\theta \right) 
\end{eqnarray}
and the convection vector components 
\begin{eqnarray}
F_{p}^{\text{C}} &=&-F\left( p,\psi ,\theta \right)   \label{Fcoll} \\
F_{\xi }^{\text{C}} &=&0
\end{eqnarray}
